\documentclass[11pt, a4paper]{article}
\usepackage[a4paper, margin=2.5cm]{geometry}
\usepackage[utf8]{inputenc}
\usepackage[T1]{fontenc}
\usepackage[norsk]{babel}
\usepackage{enumitem}
\setlist[itemize]{label=-}
\usepackage{amsmath}
\usepackage{booktabs}
\usepackage{listings}
\usepackage{hyperref}

\setlength{\parindent}{0pt}

% --- Document Info ---
\title{STK1110 - Statistiske metoder og dataanalyse: \\ Obligatorisk øving 1 av 2}
\author{Vegard Nyquist Borgeraas}
\date{Høst 2026}

\begin{document}

\maketitle

\vspace{1cm}

\section*{Oppgave 1}

\subsection*{(a)}
Ved å sette de empiriske momentene lik de teoretiske, får vi følgende likningssystem:
\begin{align*}
    M_1 &:= \frac{1}{n} \sum_{i=1}^{n} x_i \overset{!}{=} \frac{\alpha}{\gamma} \\[1ex]
    M_2 &:=\frac{1}{n} \sum_{i=1}^{n} x_i^2 \overset{!}{=} \text{Var}(X) + E[X]^2 = \frac{\alpha}{\gamma^2} + \left(\frac{\alpha}{\gamma}\right)^2
\end{align*}

Løser vi dette systemet for de to ukjente parametrene, får vi momentestimatorene:
\begin{align*}
    \hat{\gamma} &= \frac{M_1}{M_2 - M_1^2} \\
    \hat{\alpha} &= \frac{M_1^2}{M_2 - M_1^2}
\end{align*}

Resultatene samsvarer med boken, ettersom vi har at:
\begin{equation*}
    \hat{\gamma} = \frac{M_1}{M_2 - M_1^2} =  \frac{1}{\hat{\beta}_{\text{bok}}}
\end{equation*}
Dette stemmer overens med at $\beta = 1/\gamma$.

R-kode for å beregne de empiriske momentene basert på forsikringsdataene:
\vspace{0.5cm}
\begin{verbatim}
    data <- scan("forsikringskrav.txt")

    M1 <- mean(data)
    M2 <- mean(data^2)

    alpha <- M1^2 / (M2 - M1^2)
    gamma <- M1 / (M2 - M1^2)
\end{verbatim}

Dette gir at $\hat{\gamma}_{\text{mom}} = 0.696649$ og  $\hat{\alpha}_{\text{mom}} = 0.02886053$.

\subsection*{(b)}
Log-likelihood-funksjonen, $l(\alpha, \gamma)$, finner vi ved å ta den naturlige logaritmen av likelihood-funksjonen. Ved å bruke vanlige logaritmeregler gjøres produktet om til en sum:
\begin{align*}
    l(\alpha, \gamma) 
    &= \sum_{i=1}^{n} \ln( f(x_i | \alpha, \gamma) ) \\
    &= \sum_{i=1}^{n} \ln \left( \frac{\gamma^\alpha}{\Gamma(\alpha)} x_i^{\alpha-1} e^{-\gamma x_i} \right) \\
    &= \sum_{i=1}^{n} \left( \alpha \ln(\gamma) - \ln(\Gamma(\alpha)) + (\alpha - 1) \ln(x_i) - \gamma x_i \right) \\
    &= n \alpha \ln(\gamma) - n \ln(\Gamma(\alpha)) + (\alpha - 1) \sum_{i=1}^{n} \ln(x_i) - \gamma \sum_{i=1}^{n} x_i
\end{align*}

\vspace{0.5cm}
R-kode for beregning av log-likelihood:
\begin{verbatim}
    l <- function(x, alpha, gamma) {
        n <- length(x)
        n * alpha * log(gamma) - n * lgamma(alpha) + 
        (alpha - 1) * sum(log(x)) - gamma * sum(x)
    }
\end{verbatim}

Dette gir at $l(\hat{\alpha}_{\text{mom}}, \hat{\gamma}_{\text{mom}}) = -27264.51$.


\subsection*{(c)}
For å finne maksimum likelihood-estimatoren for $\gamma$ tar vi den partiellderiverte av $l(\alpha, \gamma)$ med hensyn på $\gamma$, og setter uttrykket lik null:
\begin{align*}
\frac{\partial l}{\partial \gamma} &= \frac{\partial}{\partial \gamma} \left( n \alpha \ln(\gamma) - n \ln(\Gamma(\alpha)) + (\alpha - 1) \sum_{i=1}^{n} \ln(x_i) - \gamma \sum_{i=1}^{n} x_i \right) \\
&=\frac{n \alpha}{\gamma} - \sum_{i=1}^{n} x_i
\end{align*}

For å finne maksimumspunktet, setter vi den deriverte lik null:
$$
\frac{n \alpha}{\gamma} - \sum_{i=1}^{n} x_i = 0
$$

Løser vi denne ligningen for $\gamma$, får vi maksimum likelihood-estimatoren:
$$
\hat{\gamma}_{\text{ML}} = \frac{n \alpha}{\sum_{i=1}^{n} x_i} = \frac{\alpha}{\frac{1}{n} \sum_{i=1}^{n} x_i} = \frac{\alpha}{\bar{x}}
$$
Fra oppgave (a) har vi at:
$$
\bar{x} = \frac{\alpha}{\gamma} \implies \gamma = \frac{\alpha}{\bar{x}}
$$

Dermed samsvarer de to estimatorene.


\subsection*{(d)}
R-kode for numerisk optimering:
\vspace{0.5cm}
\begin{verbatim}
    negloglikgamma <- function(logalpha,x=x) {
        n <- length(x)
        alpha <- exp(logalpha)
        gamma <- alpha/mean(x)
        logL <- n*alpha*log(gamma)-n*lgamma(alpha)+
        (alpha-1)*sum(log(x))-gamma*sum(x)
        -logL
    }

    fit.ml <- optim(log(alpha), negloglikgamma, x=x, method="BFGS")

    alpha_ml <- exp(fit.ml$par)
    gamma_ml <- alpha_ml / mean(x)

    logL_ml <- -fit.ml$value
\end{verbatim}

Dette gir at $\hat{\gamma}_{\text{ML}} = 0.05751589$ og $\hat{\alpha}_{\text{ML}} = 1.388346$, og dermed at $l(\hat{\alpha}_{\text{ML}}, \hat{\gamma}_{\text{ML}}) = -26488.28$. Log-likelihood-verdien for ML-estimatene er helt rimelig fordi den er høyere enn verdien fra momentestimatene i (b), slik man alltid forventer når man maksimerer funksjonen.

\subsection*{(e)}
R-kode for bootstrap-simulering:
\begin{verbatim}
    B <- 1000 # Antall bootstrap-simuleringer
    n <- length(x)

    alpha_boot <- numeric(B)
    gamma_boot <- numeric(B)

    for (i in 1:B) {
        x_b <- sample(x, size = n, replace = TRUE)
        
        fit_b <- optim(log(alpha), negloglikgamma, x = x_b, method = "BFGS")
        
        alpha_boot[i] <- exp(fit_b$par)
        gamma_boot[i] <- alpha_boot[i] / mean(x_b)
    }
\end{verbatim}

Da får vi følgende anslag for standardfeilen (SE) til maksimum likelihood-estimatene:
\begin{align*}
    \text{SE}(\hat{\alpha}_{\text{ML}}) &\approx 0.0297 \\
    \text{SE}(\hat{\gamma}_{\text{ML}}) &\approx 0.0019
\end{align*}

Og vi får følgende 95 \% konfidensintervaller:
\begin{align*}
    \text{95 \% KI for } \alpha &: \quad [1.335, \, 1.450] \\
    \text{95 \% KI for } \gamma &: \quad [0.0542, \, 0.0614]
\end{align*}


\subsection*{(f)}
R-kode for konfidensintervaller for $\mu$:
\begin{verbatim}
    mu_boot <- alpha_boot / gamma_boot

    ci_mu_95 <- quantile(mu_boot, c(0.025, 0.975))
    ci_mu_99 <- quantile(mu_boot, c(0.005, 0.995))
\end{verbatim}

Da får vi følgende konfidensintervaller for forventningsverdien $\mu$:
\begin{align*}
    \text{95 \% KI for } \mu &: \quad [23.43, \, 24.81] \\
    \text{99 \% KI for } \mu &: \quad [23.29, \, 24.96]
\end{align*}

Vi tester hypotesene $H_0: \mu = 25$ mot $H_a: \mu \neq 25$. 

Siden $\mu = 25$ ligger utenfor det brede 99 \% konfidensintervallet, forkaster vi nullhypotesen på signifikansnivå $\alpha = 0.01$. Ettersom 99 \%-intervallet inneholder 95 \%-intervallet, forkastes nullhypotesen i tillegg på signifikansnivå $\alpha = 0.05$.

Da får vi følgende konklusjon for P-verdien:
\begin{align*}
    \text{P-verdi} &< 0.01
\end{align*}

\newpage
\section*{Oppgave 2}

\subsection*{a)}

Siden den sanne variansen $\sigma^2$ er ukjent og vi antar at målingene er uavhengige og normalfordelte, benytter vi en t-fordeling med $n-1 = 14$ frihetsgrader. 

R-kode for utregning av konfidensintervall:
\begin{verbatim}
    x <- c(525, 587, ..., 562)
    n <- length(x)

    x_bar <- mean(x)
    s <- sd(x)
    t_val <- qt(0.975, df = n - 1)

    nedre <- x_bar - t_val * s / sqrt(n)
    ovre <- x_bar + t_val * s / sqrt(n)
\end{verbatim}

Datasettet gir følgende konfidensintervall for forventet antocyaninnhold, $\mu$:
\begin{align*}
    \text{95 \% KI for } \mu &= \bar{x} \pm t_{0.025, 14} \frac{s}{\sqrt{n}} \\
    &= [543.85, \, 575.48]
\end{align*}

\subsection*{b)}

R-kode for simulering av dekningsgrad: 
\begin{verbatim}
    B <- 10000
    n <- 15
    mu_sann <- 558
    sigma_sann <- 30

    sim_data <- matrix(
        rnorm(B * n, mean = mu_sann, sd = sigma_sann), 
        nrow = B, 
        ncol = n
    )

    x_bar <- rowMeans(sim_data)
    s <- apply(sim_data, 1, sd)
    t_val <- qt(0.975, df = n - 1)
    feilmargin <- t_val * s / sqrt(n)

    nedre_grense <- x_bar - feilmargin
    ovre_grense <- x_bar + feilmargin

    dekker_mu <- (nedre_grense <= mu_sann) & (ovre_grense >= mu_sann)

    andel <- mean(dekker_mu)
\end{verbatim}

Andelen ble $0.9517$, noe som er en god tilnærming til 95 \%. Dette illustrerer definisjonen av et konfidensintervall: at omtrent 95 \% av alle slike konstruerte intervaller faktisk vil inneholde den sanne forventningsverdien $\mu$.


\subsection*{c)}
R-kode for tilnærmet konfidensintervall:
\begin{verbatim}
    feilmargin_stor <- 1.96 * s / sqrt(n)

    nedre_grense_stor <- x_bar - feilmargin_stor
    ovre_grense_stor <- x_bar + feilmargin_stor

    dekker_mu_stor <- (nedre_grense_stor <= mu_sann) & (ovre_grense_stor >= mu_sann)

    andel_stor <- mean(dekker_mu_stor)
\end{verbatim}

Når vi bruker det tilnærmede intervallet, viser simuleringen at andelen synker til $0.9325$. Dette stemmer med at denne tilnærmingen antar at vi har mer informasjon om standardavviket enn vi faktisk har, og dermed lager et for smalt intervall.

\subsection*{d)}
For å konstruere et konfidensintervall for standardavviket $\sigma$ i en normalfordeling, bruker vi kji-kvadrat-fordelingen. Vi vet at størrelsen $\frac{(n-1)S^2}{\sigma^2}$ er $\chi^2$-fordelt med $n-1$ frihetsgrader. Ved å snu på dette uttrykket, får vi følgende formel for et 95 \% konfidensintervall for $\sigma$:
\begin{equation*}
    \left[ \sqrt{\frac{(n-1)s^2}{\chi^2_{0.975, 14}}}, \, \sqrt{\frac{(n-1)s^2}{\chi^2_{0.025, 14}}} \right]
\end{equation*}

\begin{verbatim}
    chi_nedre <- qchisq(0.025, df = n - 1)
    chi_ovre  <- qchisq(0.975, df = n - 1)

    nedre_grense_sigma <- sqrt((n - 1) * s^2 / chi_ovre)
    ovre_grense_sigma  <- sqrt((n - 1) * s^2 / chi_nedre)

    dekker_sigma <- (nedre_grense_sigma <= sigma_sann) & (ovre_grense_sigma >= sigma_sann)
    andel_sigma <- mean(dekker_sigma)
\end{verbatim}

Når vi kjører denne simuleringen finner vi at andelen er $0.9489$, som er tilnærmet $95$ \%, akkurat som forventet.

\subsection*{e)}
R-kode for simulering med t-fordelte data: 
\begin{verbatim}
    sim_data <- mu_sann + sigma_sann * matrix(
        rt(B * n, df=7), 
        nrow = B, 
        ncol = n
    )
\end{verbatim}

Når vi kjører simuleringen er andelen $0.9511$, som fremdeles ligger svært nær $95$ \%. 

Dette resultatet viser at metoden for å lage konfidensintervall for forventningsverdien er svært robust mot avvik i antakelsen om normalfordeling.

\subsection*{f)}
R-kode for konfidensintervall for standardavviket med t-fordelte data: 
\begin{verbatim}
    sigma_tilde <- sigma_sann * sqrt(7 / (7 - 2))
    dekker_sigma <- (nedre_grense_sigma <= sigma_sann) & 
                    (ovre_grense_sigma >= sigma_sann)
    andel_sigma_tilde <- mean(dekker_sigma_tilde)
\end{verbatim}

Andelen ble $0.8823$, som er betydelig lavere enn 95 \%. Dette viser at metoden for å lage konfidensintervall for standardavviket er mye mer sensitiv for avvik fra antakelsen om normalfordeling.

\end{document}